\documentclass[11pt]{article}
\usepackage[margin=1in]{geometry}
\usepackage{amsmath,amssymb,amsthm,mathtools}
\usepackage[hidelinks,pdfusetitle]{hyperref}
\usepackage{microtype}
\usepackage{enumitem}
\usepackage{booktabs}
\usepackage{array}
\newcommand{\DOI}[1]{\mbox{\href{https://doi.org/#1}{\nolinkurl{#1}}}}
\newcommand{\kb}{k_{\mathrm B}}
\newcommand{\Tlin}{T_{\mathrm{lin}}}
\newcommand{\Om}{\Omega}
\newcommand{\OmE}{\Omega_E}
\newcommand{\CE}{C_E}
\newcommand{\Fcal}{\mathcal F}
\newcommand{\muS}{\mu^{*}}
\newtheorem{theorem}{Theorem}[section]
\newtheorem{lemma}[theorem]{Lemma}
\newtheorem{corollary}[theorem]{Corollary}
\newtheorem{proposition}[theorem]{Proposition}
\theoremstyle{definition}
\newtheorem{definition}[theorem]{Definition}
\newtheorem{remark}[theorem]{Remark}
\theoremstyle{remark}
\newcommand{\proved}{\textnormal{\textsc{(proved)}}}
\newcommand{\certified}{\textnormal{\textsc{(certified)}}}
\newcommand{\numerical}{\textnormal{\textsc{(numerical)}}}
\allowdisplaybreaks
\setlength{\emergencystretch}{3em}

\title{The Coulomb pseudopotential and the optimality of the Einstein spectrum\\ for the linear Eliashberg ceiling}
\author{Davis Rens\\SuperMatics Inc.\\\href{https://orcid.org/0000-0001-5839-5197}{ORCID 0000-0001-5839-5197}}
\date{20 November 2026}
\begin{document}
\maketitle

\begin{abstract}
We ask whether the single Einstein line still maximizes the linear pairing onset of
isotropic Migdal--Eliashberg theory when a pairing-only Coulomb repulsion acts through a
cutoff that is \emph{common} to every spectrum compared at fixed coupling $\lambda$ and
spectral moment $M$. At zero repulsion the answer is yes, and the maximizing value is the
sharp ceiling $\kb\Tlin\le C(\lambda)\sqrt M$, whose large-coupling limit is
$C_\infty=0.18272624777232952$~\cite{univglue,anisosharp}. We show first that the ceiling value is Coulomb robust:
the repulsion enters the folded stability form as a positive semidefinite rank-one term,
so no spectrum at any $\muS\ge0$ exceeds $C(\lambda)\sqrt M$. The optimizer is not robust.
Using exact rational arithmetic and two independently written verifiers we certify three
two-line spectra that share $(\lambda,M,\muS,\omega_c)$ with the Einstein line and have a
strictly larger linear onset, at $\muS/\lambda=2/3$, $3/4$ and $1$, with temperature
gains of at least $2.59$, $29.09$ and $82.57$ percent. We record seven certified Einstein
brackets, including the zero-repulsion calibration
$C(1)\in[0.11461945,0.11461957]$. A separate section goes below one half: four further
counterexamples, certified by a second pair of independent exact verifiers in the same two
roles, confined to $\muS/\lambda=0.4786$ to $0.4890$ at couplings $\lambda=25.6$ to $28$,
so that $\muS/\lambda=1/2$ is not a threshold of the model, together with proved
optimality in three regimes, a two-line reduction that makes every comparison
two-dimensional, and floating-point data locating the failure. Statements carry the
labels
\textsc{proved}, \textsc{certified}, \textsc{numerical} and \textsc{open}, which refer to
mathematical or computational evidence and never to measurements. The bounds are on the
model quantity $\Tlin$; no measured transition temperature is discussed.
\end{abstract}

\section{The question and its position}

For a fixed coupling $\lambda$ and a fixed second spectral moment $M$ of the
electron-phonon glue, linear Migdal--Eliashberg theory has a ceiling: the linear onset
temperature $\Tlin$ obeys $\kb\Tlin\le C(\lambda)\sqrt M$, the coefficient $C(\lambda)$
increases with $\lambda$, and the bound is sharp at $C_\infty=0.18272624777232952$~\cite{univglue,anisosharp} as
$\lambda\to\infty$. The spectrum that attains the fixed-$(\lambda,M)$ value is a single
Einstein line. These statements are prior work: transfer-by-transfer comparison,
temperature monotonicity and strict single-line optimality~\cite{rens}, the same sharp
spectral-moment bound at $\muS=0$~\cite{kondratev}, the certification of the
$\gamma=2$ constant~\cite{gamma} as used in the sharp scalar
ceiling~\cite{finiteband,gsi}, and the multipoint moment-norm ceilings for coupled
branches~\cite{univglue,optimum,anisosharp,capstone}. The classical strong-coupling
expressions~\cite{mcmillan,allendynes} are the formulas in which one Einstein line stands
for the whole glue spectrum. Moussa and Cohen identified the largest moment eigenvalue
rather than the angular average as the flat multiband strong-coupling
scale~\cite{moussa}, and Kiessling, Altshuler and Yuzbashyan formulated the flat scalar
spectral problem on which the ceiling rests~\cite{kay}. This paper follows the notation of
those releases and is a companion to them~\cite{univglue,optimum,anisosharp,capstone,finiteband,etapair}.

Every one of those results is a zero-Coulomb statement. In a phonon-mediated
superconductor the Coulomb repulsion is not zero; what enters the pairing channel is the
renormalized pseudopotential $\muS=\mu/(1+\mu\ln(E_F/\omega_c))$ that results when a bare
Coulomb coupling $\mu$ is combined with the cutoff $\omega_c$, and that renormalization
maps a large bare $\mu$ to a modest $\muS$ without reference to $\lambda$. How accurate
that renormalization is, is not at issue here. The question of this paper is what a
nonzero pairing-only $\muS$ does to Einstein optimality when the cutoff $\omega_c$ is a
single number shared by every spectrum compared at the same $(\lambda,M)$.

The answer has two halves, and they point in opposite directions.
\begin{enumerate}[label=(\arabic*),leftmargin=2.2em]
\item \emph{The ceiling value survives.} The repulsion enters the folded stability form
as a positive semidefinite rank-one term; hence no spectrum at any $\muS\ge0$ has an
onset above the zero-repulsion Einstein value $C(\lambda)\sqrt M$
(Corollary~\ref{cor:value}, \textsc{proved}).
\item \emph{The optimizer does not.} Above $\muS/\lambda=1/2$ there are two-line spectra
with the same $(\lambda,M,\muS,\omega_c)$ as Einstein whose linear onset is strictly
larger. Three such instances, at $\muS/\lambda=2/3$, $3/4$ and $1$, are certified in
Section~\ref{sec:cert} with exact rational witnesses and two independent verifiers, with
gains of at least $2.59$, $29.09$ and $82.57$ percent in temperature. Below one half the
failure persists but only in a narrow window of extreme coupling: Section~\ref{sec:half}
certifies four further instances down to $\muS/\lambda=0.4786$, all of them at
$\lambda=25.6$ to $28$, and proves that Einstein is optimal in three regimes.
\end{enumerate}
So the sharp value of the ceiling is not at stake; the identity of the maximizing spectrum
is. A user of the ceiling who assumes that the single line is the worst case at fixed
$(\lambda,M)$ is safe in value and unsafe in spectrum once the pseudopotential is large
relative to the coupling.

Section~\ref{sec:model} fixes the model, Section~\ref{sec:reduce} the reduction to three
parameters, Section~\ref{sec:dom} the Einstein dominance that holds at $\muS=0$ and its
failure at $\muS>0$, Section~\ref{sec:cert} the certified counterexamples and the
certification method, Section~\ref{sec:brackets} the certified brackets and the
zero-repulsion calibration, Section~\ref{sec:num} the direct-discretization cross-check,
Section~\ref{sec:half} the counterexamples below one half with the regimes in which
Einstein is proved optimal there, Section~\ref{sec:not} what is not claimed, and
Section~\ref{sec:open} the open problems.

Throughout, a statement labelled \textsc{proved} has a proof in the text;
\textsc{certified} means that a finite exact-arithmetic certificate has been checked by
two independently written verifiers; \textsc{numerical} means floating-point evidence
only; \textsc{open} means not established in either direction.

\section{Model and the common cutoff}\label{sec:model}

We use the isotropic, flat-density-of-states, infinite-band linearized Migdal--Eliashberg
model with an even gap in the conventions of the companion releases
\cite{finiteband,univglue}. Energies are in absolute units and the Matsubara frequencies
are $\omega_n=(2n+1)\pi\kb T$. The glue defines
\begin{equation}\label{eq:const}
\lambda=2\int_0^\infty\frac{\alpha^2F(\omega)}{\omega}\,d\omega,\qquad
M=2\int_0^\infty\alpha^2F(\omega)\,\omega\,d\omega,\qquad
\OmE=\sqrt{M/\lambda},
\end{equation}
the probability measure $dp=2\alpha^2F(d\omega)/(\lambda\omega)$ on
$u=\omega^2/\OmE^2$ has $\int dp=\int u\,dp=1$, and the transfer weights at temperature
$T$ are
\begin{equation}\label{eq:transfers}
a_r(q)=\lambda\int\frac{u}{u+qr^2}\,dp,\qquad q=\Big(\frac{2\pi\kb T}{\OmE}\Big)^2,
\end{equation}
so that $a_0=\lambda$, $a_r>0$ and $\sum_r a_r<\infty$. A spectrum is any probability
measure $p$ with $\int u\,dp=1$; the Einstein line in these variables is
$p=\delta_1$, and the two-line family is $p=(1-w)\delta_{u_1}+w\delta_{u_2}$ with
$u_1=(1-wu_2)/(1-w)$ fixed by the moment condition.

The Coulomb repulsion is \emph{pairing only}: it acts as $\muS c_nc_m$ with the strict
mask $c_n=\mathbf 1_{|\omega_n|<\omega_c}$ on both indices, and it never enters the normal
self-energy $Z$. This is the standard pseudopotential truncation; keeping $\muS$ out of
$Z$ is what makes the problem a competition between two attractive kernels rather than a
change of the normal state. The cutoff is \emph{common}: one number $\omega_c$ per
instance, identical for every spectrum compared there. It is natural to measure it
against the Einstein frequency, so we write
\begin{equation}\label{eq:R}
R=\omega_c/\OmE=\omega_c\sqrt{\lambda/M},\qquad \omega_c=R\OmE,
\end{equation}
and the pseudopotential is measured against the coupling, by the ratio
$\muS/\lambda$.

With the normal diagonal $D_n=Z_n|\omega_n|/(\pi\kb T)=|2n+1|+a_0+2\sum_{j=1}^{i(n)}a_j$ and the
pairing kernel $K_{nm}=a_{|n-m|}-\muS c_nc_m$, the linearized gap equation is the
eigenproblem $\lambda_{\max}(D^{-1/2}KD^{-1/2})=1$. We work with the equivalent folded
stability form on the even subspace $h_n=h_{-n-1}$, indexed by
$i=\max(n,-n-1)\ge0$ with $d_i=|2i+1|$. For $A_i=\sum_{j=1}^{i}a_j$ the form is
\begin{equation}\label{eq:form}
\Fcal(h)=\sum_{i\ge0}\big(d_i+2A_i\big)h_i^2
-2\sum_{i<j}a_{j-i}h_ih_j
-\sum_{i,j\ge0}a_{i+j+1}h_ih_j
+2\muS\Big(\sum_{i<k}h_i\Big)^2 ,
\end{equation}
where $k=k(q,R)$ is the number of folds inside the mask,
\begin{equation}\label{eq:mask}
i<k\iff (2i+1)^2q<4R^2 ,
\end{equation}
and the quotient identity
\begin{equation}\label{eq:quotient}
\frac{h^{\rm T}Kh}{h^{\rm T}Dh}
=1-\frac{\Fcal(h)}{\sum_{i\ge0}\big(d_i+a_0+2A_i\big)h_i^2}
\end{equation}
holds for finitely supported $h$, with a strictly positive denominator. The medium is
stable if and only if $\Fcal\succeq0$, that is, if and only if
$\lambda_{\max}(D^{-1/2}KD^{-1/2})\le1$, and the linear onset is
$\kb\Tlin=\sup\{\kb T:\Fcal\not\succeq0\}$. It is the supremum of temperatures at which the
linearized form has a negative direction, not a root of a nonlinear gap equation, and it
is not the transition temperature of any material. The scale-free coefficient carried by
a given onset is
\begin{equation}\label{eq:C}
C=\frac{\kb\Tlin}{\sqrt M}=\frac{\sqrt q}{2\pi\sqrt\lambda}.
\end{equation}

Two facts about \eqref{eq:form} are used constantly. First, the Coulomb term is
$2\muS(\sum_{i<k}h_i)^2$, a positive semidefinite rank-one form supported on the masked
block, so it can only raise the form. Second, the zero-repulsion part of \eqref{eq:form}
is the object whose optimality was settled in~\cite{rens,kondratev}: its onset at fixed
$(\lambda,M)$ is largest for $p=\delta_1$, and we write $q_{\rm ph}(\lambda)$ for that
Einstein onset at $\muS=0$, with $\kb T_{\rm ph}=\sqrt{q_{\rm ph}}\,\OmE/(2\pi)$.

\section{Structural reduction at fixed coupling and moment}\label{sec:reduce}

\begin{theorem}[Reduction to three parameters]\label{thm:reduce}\proved
Fix $0<\lambda,M<\infty$ and $\muS\ge0$. Rescaling all frequencies by $\OmE$ maps every
instance to the triple $(\lambda,R^2,\muS)$ with $R=\omega_c/\OmE$; the mask is decided by
\eqref{eq:mask}; and the onset, the eigenproblem and the stability form depend on the
spectrum only through the weights $a_r(q)$ of \eqref{eq:transfers}. Consequently two
spectra with equal $\lambda$, equal $M$ and equal $\muS$ compete on the same footing at
equal $R$ and at no other $R$.
\end{theorem}
\begin{proof}
$\OmE$ is determined by $(\lambda,M)$, so all quantities in \eqref{eq:transfers} and
\eqref{eq:mask} are functions of the triple. The eigenproblem and \eqref{eq:form} are
scale invariant once $q$ and $R^2$ are used.
\end{proof}

Theorem~\ref{thm:reduce} is what makes the question well posed. A cutoff fixed as a
multiple of $\OmE$ is common by construction, and Einstein's own mask then has the same
width $R$ as everyone else's. The distinction matters, because a cutoff that follows the
highest frequency of each spectrum instead has a different and degenerate
answer~\cite{rens,finiteband}. In the fixtures used below every mixed spectrum satisfies
$\sum_i\lambda_i=\lambda$ and $\sum_i\lambda_iu_i=\lambda$, so $M$ and $\OmE$ coincide with
Einstein's and $R^2$ is literally the same rational number.

The zero-repulsion ceiling, in the notation of \eqref{eq:C}, reads
\begin{equation}\label{eq:ceiling}
\kb\Tlin\le C(\lambda)\sqrt M,\qquad C(\lambda)\Uparrow C_\infty=0.18272624777232952,
\end{equation}
with $C(\lambda)$ the value \eqref{eq:C} at the Einstein onset and $C_\infty$ the
certified $\gamma=2$ constant~\cite{gamma,finiteband}. We keep \eqref{eq:ceiling} as the
reference throughout: nothing below changes it.

\section{Einstein dominance and its failure}\label{sec:dom}

\subsection{What holds at zero repulsion}

\begin{lemma}[Jensen dominance of the Einstein transfers]\label{lem:jensen}\proved
For every spectrum, every $q>0$ and every $r\ge1$,
\begin{equation}\label{eq:jensen}
a_r(q)=\lambda\int\frac{u}{u+qr^2}\,dp
\le\lambda\frac{\int u\,dp}{\int u\,dp+qr^2}=\frac{\lambda}{1+qr^2}=a_r^{\rm E}(q),
\end{equation}
with equality for some $r\ge1$ if and only if $p=\delta_1$.
\end{lemma}
\begin{proof}
The map $u\mapsto u/(u+qr^2)$ is strictly concave on $(0,\infty)$ and $\int u\,dp=1$.
Jensen gives the inequality, and strict concavity with a non-degenerate spectrum gives
strictness.
\end{proof}

Lemma~\ref{lem:jensen} is the analytic content of the zero-repulsion ceiling: every
transfer of every spectrum is dominated by the corresponding Einstein transfer, entry by
entry. Its combination with an exact rewriting of \eqref{eq:form} gives the fixed-moment
optimum~\cite{rens}.

\begin{lemma}[Folded Laplacian identity]\label{lem:laplace}\proved
With $S_1=\sum_{r\ge1}a_r$ and $T_i=\sum_{r>i}a_r=S_1-A_i$, the form \eqref{eq:form} at
$\muS=0$ equals
\begin{equation}\label{eq:display}
\Fcal_0(h)=\sum_{i\ge0}(2i+1)h_i^2-2\sum_{i\ge0}h_i^2T_i
+\tfrac12\sum_{i,j\ge0}a_{|i-j|}(h_i-h_j)^2
+\tfrac12\sum_{i,j\ge0}a_{i+j+1}(h_i-h_j)^2 ,
\end{equation}
for finitely supported $h$, and the two expressions agree identically.
\end{lemma}
\begin{proof}
Use $\sum_{i<j}a_{j-i}(h_i-h_j)^2=S_1\sum_ih_i^2+\sum_jA_jh_j^2-2\sum_{i<j}a_{j-i}h_ih_j$
and $\sum_{i,j}a_{i+j+1}(h_i-h_j)^2=2\sum_ih_i^2T_i-2\sum_{i,j}a_{i+j+1}h_ih_j$; add the
result to the diagonal $\sum_i(d_i+2A_i)h_i^2$. Three bookkeeping points matter: both
Laplacian sums carry the factor $1/2$ because each unordered pair is counted twice; the
diagonal carries $-2\sum_{r>i}a_r$, not $-\sum_{r>i}a_r$; and $a_0$ cancels between the
$Z$ diagonal and the kernel diagonal, so neither \eqref{eq:form} nor \eqref{eq:display}
shows it. The identity was also verified in exact rational arithmetic on 120 random
trials with mixed-sign and all-positive $h$, random two-line spectra, random $\muS$ and a
random mask, and the mutant display that drops the tail term was detected in all 120.
\end{proof}

The display \eqref{eq:display} is the form a reader should keep in mind: a kinetic
diagonal, two Laplacian couplings to the transfer tails, and nothing else. The
$a_{i+j+1}$ sum is the even-frequency image of the negative Matsubara axis. It is not
optional, and it is not a source of unexplained error: any discretization that spans the
negative Matsubara axis carries it automatically, and it is lost only when that axis is
dropped from the start.

\subsection{The Coulomb block and why dominance does not transfer}

\begin{proposition}[Rank-one structure at positive $\muS$]\label{prop:rank}\proved
For every spectrum, every $\muS\ge0$ and every $q$,
\begin{equation}\label{eq:rank}
\Fcal(h)=\Fcal_0(h)+2\muS\Big(\sum_{i<k}h_i\Big)^2 .
\end{equation}
The repulsion raises the form on the masked block and on the masked block only; the
masked block is the lowest $k$ folds of \eqref{eq:mask}, and outside it the form is
untouched.
\end{proposition}
\begin{proof}
Immediate from \eqref{eq:form}, since $c_n$ enters only through the pairing kernel and the
normal diagonal is $\muS$-independent.
\end{proof}

\begin{corollary}[The ceiling value is Coulomb robust]\label{cor:value}\proved
Let $\muS\ge0$ and fix $(\lambda,M)$. Every spectrum has
$\kb\Tlin\le C(\lambda)\sqrt M\le C_\infty\sqrt M$. In particular $q_{\rm ph}$
is a hard upper bound for the onset of every spectrum at every $\muS\ge0$.
\end{corollary}
\begin{proof}
By \eqref{eq:rank}, $\Fcal\succeq\Fcal_0$ at $\muS\ge0$, so a temperature at which the
repulsive form has a negative direction is a temperature at which the zero-repulsion form
has one. The zero-repulsion onset of every spectrum is at most the Einstein
value~\cite{rens,kondratev}, which gives the first inequality; the second is
\eqref{eq:ceiling}. Also, for $q\ge q_{\rm ph}(\lambda)$ the zero-repulsion form is
positive semidefinite for every spectrum, and the rank-one term cannot destroy that.
\end{proof}

Corollary~\ref{cor:value} is the reassuring half of the paper, and it is the reason the
counterexamples below are not a weakening of any ceiling. What fails at positive $\muS$ is
the passage from Lemma~\ref{lem:jensen} to the forms. Entrywise dominance of the transfers
does not imply dominance of the quadratic forms, because the kernel enters with both
signs: in \eqref{eq:display} a transfer $a_r$ appears as a reward through $-2T_i$ on the
diagonal and as a penalty through the two Laplacians. The obstruction is visible in single
entries of \eqref{eq:form}: there the diagonal is
$F_{ii}=d_i+2A_i-a_{2i+1}+2\muS\,[i<k]$, which \emph{rises} when transfers are added to
the near zone, since $A_i$ sums them while $a_{2i+1}$ is subtracted, and every off-diagonal
entry $F_{ij}=-a_{|i-j|}-a_{i+j+1}$ ($i\ne j$) \emph{falls}. The same transfer is thus a
reward on the diagonal and a penalty off it, so a statement that all transfers grow fixes
no ordering of the forms. At $\muS=0$ the deficit does
not matter, because the Einstein line still wins the comparison by other means. At
$\muS>0$ it does matter, for the following reason.

\begin{remark}[The mechanism]\label{rem:mech}
At $\muS=0$ a non-Einstein spectrum pays twice: moving weight out of the near zone lowers
every transfer, in particular $a_1$, and $a_1$ is what pays for the normal mass at small
$j$. A two-line spectrum can try to dodge the payment by putting a small weight $w\ll1$ at
a large $u_2\gg q$: the moment condition then forces $u_1\approx1/(1-w)<1$, so the bulk of
the weight sits just below the Einstein line and $a_1$ stays close to
$\lambda/(1+q)$. That is still a loss at $\muS=0$. At $\muS>0$ the mask removes the lowest
folds from the competition, and the repulsion penalizes a form in proportion to
$(\sum_{i<k}h_i)^2$. The Einstein instability at its crossing is spread over the whole
near zone, so it collects the whole penalty; the mixture's instability is carried by a
narrow set of folds just outside the mask, which collects much less. Beyond a threshold
ratio $\muS/\lambda$ the trade reverses and the mixture becomes the more unstable
spectrum. Section~\ref{sec:cert} exhibits the reversal as an exact certificate for three
choices of the ratio.
\end{remark}

A direct measurement of the same competition, in the full Matsubara line at the wide
counterexample's witness $q=1/200$, gives for the largest eigenvalue of
$D^{-1/2}KD^{-1/2}$ the pair $4.4033\to0.7556$ for the Einstein line against
$3.2054\to1.0975$ for the mixture, the two numbers of each pair being $\muS=0$ and
$\muS=3/2$ (\textsc{numerical}). The repulsion removes the Einstein instability far more
effectively than the mixture's.

\subsection{Two regimes where the repulsion cannot disturb the optimum}

The failure described in the next section needs both a mask wide enough to compete with
the near zone and a ratio $\muS/\lambda$ large enough to make the rank-one term bite. Each
requirement can be removed on its own, and both removals are one-line computations, so
they are recorded here as observations rather than as separate results.

\emph{Empty mask at and above the crossing.} By \eqref{eq:mask} the mask contains the fold
$i=0$ exactly when $q<4R^2$. If $4R^2<q_{\rm ph}(\lambda)$, then on $[q_{\rm ph},\infty)$
the mask is empty, so $\Fcal=\Fcal_0$ there and every spectrum is stable, the zero-repulsion
form being positive semidefinite at and above $q_{\rm ph}$ for every spectrum. By
Corollary~\ref{cor:value} no spectrum has an onset above its own zero-repulsion onset,
which is strictly below $q_{\rm ph}$ for every non-Einstein spectrum, while on the nonempty
open interval $(4R^2,q_{\rm ph}(\lambda))$ the mask is still empty, so $\Fcal=\Fcal_0$ and
the Einstein form remains indefinite arbitrarily close to $q_{\rm ph}(\lambda)$.
Hence at such $R$ the Einstein line strictly maximizes the onset for every $\muS\ge0$, and
the counterexamples of Section~\ref{sec:cert} require a mask that is still alive at the
crossing.

\emph{Parameter-free stability at large $q$.} For a spectrum whose transfers are written
in the Laplacian form of \eqref{eq:form}, one has $u/(u+qr^2)\le u/(qr^2)$, hence
$a_r\le\lambda/(qr^2)$ for $r\ge1$, hence $S_1=\sum_{r\ge1}a_r\le\pi^2\lambda/(6q)$, which
is at most $\pi^2/24$ once $q\ge4\lambda$. In the equivalent transfer decomposition
\begin{equation}\label{eq:decomp}
\begin{aligned}
\Fcal(h)&=\sum_{i\ge0}(2i+1)h_i^2+2\muS\Big(\sum_{i<k}h_i\Big)^2+\sum_{j\ge1}a_jL_j(h),\\
L_j(h)&=-\sum_{i=0}^{j-1}h_ih_{j-1-i}-2\sum_{i\ge0}h_{i+j}(h_i-h_{i+j}).
\end{aligned}
\end{equation}
the three pieces of $\sum_ja_jL_j$ are $-\sum_{i,l}a_{i+l+1}h_ih_l$,
$-2\sum_{i<l}a_{l-i}h_ih_l$ and $2\sum_iA_ih_i^2$, which together with the diagonal
$(2i+1)h_i^2$ is exactly \eqref{eq:form}. Since
$-\sum_{i<j}h_ih_{j-1-i}\ge-\sum_{i<j}h_i^2$ and $-2h_ih_{i+j}\ge-h_i^2-h_{i+j}^2$, the
difference term satisfies $L_j(h)\ge-2\sum_{i\ge0}h_i^2$ for every $j\ge1$ and every real
$h$, so
\begin{equation}\label{eq:asym}
\Fcal\succeq\Big(1-2S_1\Big)I\succeq\Big(1-\frac{\pi^2}{12}\Big)I\succ\tfrac16 I
\qquad\text{for every spectrum when } q\ge4\lambda .
\end{equation}
Above $q=4\lambda$ therefore no spectrum is unstable, for any mask and any $\muS$, with no
certificate needed. Neither observation reaches the interior between a mask that dies at
the crossing and a coupling that stabilizes the form on its own, and the counterexamples
live in that interior.

\section{Certified counterexamples above one half}\label{sec:cert}

\begin{theorem}[Three certified failures of Einstein optimality]\label{thm:cert}\certified
For each of the three parameter sets of Table~\ref{tab:cex} there exists a rational
two-line spectrum with the same $(\lambda,M,\muS,\omega_c)$ as the Einstein line, hence
the same $\OmE=\sqrt{M/\lambda}$ and the same $R$, whose linear onset is strictly larger:
$\kb\Tlin^{\rm mix}\ge\sqrt{q_{\rm mix}}\,\OmE/(2\pi)$ with $q_{\rm mix}>q_{\rm hi}$, where
$[q_{\rm lo},q_{\rm hi}]$ is the certified Einstein bracket of that row.
\end{theorem}

\begin{table*}[t]
\centering
\scriptsize
\setlength{\tabcolsep}{3pt}
\renewcommand{\arraystretch}{1.15}
\begin{tabular}{@{}lrrccccccc@{}}
\toprule
case & $\lambda$ & $\muS$ & $\muS/\lambda$ & $R^2$ & $q_E$ bracket & $\CE$ bracket &
spectrum & $q_{\rm mix}$ & gain\\
\midrule
E1 & $3$ & $2$ & $2/3$ & $9/4$ & $[0.0083,\,0.0095]$ & $[0.00837140,\,0.00895615]$ &
$0.984975$ at $0.1$, $0.015025$ at $60$ & $1/100$ & $\ge2.59\%$\\
E2 & $2$ & $3/2$ & $3/4$ & $4$ & $[0.0009,\,0.003]$ & $[0.00337618,\,0.00616405]$ &
$0.995398$ at $0.08$, $0.004602$ at $200$ & $1/200$ & $\ge29.09\%$\\
E3 & $4$ & $4$ & $1$ & $4$ & $[0.000258,\,0.003]$ & $[0.00127820,\,0.00435864]$ &
$0.990793$ at $0.08$, $0.009207$ at $100$ & $1/100$ & $\ge82.57\%$\\
\bottomrule
\end{tabular}
\caption{The three certified counterexamples. The Einstein bracket is exact rational and
outward rounded here; $C_E=\sqrt{q}/(2\pi\sqrt\lambda)$ is the Einstein coefficient of
\eqref{eq:C}. The two-line spectra are shown to six decimals; the certified objects are
the exact rationals $590/599$ at $u=1/10$ together with $9/599$ at $u=60$ (E1),
$4975/4998$ at $u=2/25$ with $23/4998$ at $u=200$ (E2), and $2475/2498$ at $u=2/25$ with
$23/2498$ at $u=100$ (E3). The gain is
$\sqrt{q_{\rm mix}/q_{\rm hi}}-1$, a lower bound on the temperature gain. In every row
$\sum_ip_i=1$ and $\sum_ip_iu_i=1$ exactly, so $\lambda$ and $M$ coincide with
Einstein's.}
\label{tab:cex}
\end{table*}

\begin{proof}[Certification]
Each row is established by a chain of finite exact statements, one for instability, one
for stability, and one for the comparison.
\emph{Cold side (instability, no truncation).} At the rational $q_{\rm lo}$ an integer
witness vector $x$ supported on the first $n$ folds is given, and the value of the
infinite-model form at $x$ is \emph{exactly} the finite quadratic form
$\sum_{i,j<n}x_iF_{ij}x_j$: every pair with $j\ge n$ multiplies $x_j=0$, and all weights
$a_r$ needed have $r\le2n-1$ and are computed exactly. No tail bound enters. A negative
value certifies $\Fcal(x)<0$ in the infinite model.
\emph{Hot side (stability on an interval).} On $[q_{\rm hi},q_{\rm ph}(\lambda))$ the
positivity of the infinite form is certified by splitting the interval at the exact mask
jump points, enclosing the head block entrywise by an integer centre $C$ and radius
$\mathcal R$ with $C-\mathcal R\le 2^{17}F_{ij}\le C+\mathcal R$ for every off-diagonal
entry, subtracting the radius as a weighted diagonal penalty (valid for arbitrary real
vectors), eliminating the whole infinite tail by a Schur complement against a uniform
Gershgorin gap of the tail, and deciding the resulting integer matrix by an exact residual
with strict diagonal dominance. Beyond $q_{\rm ph}$ no work is needed by
Corollary~\ref{cor:value}. \emph{Comparison.} $q_{\rm mix}>q_{\rm hi}$ together with the
cold witness at $q_{\rm lo}$ brackets the Einstein onset inside $[q_{\rm lo},q_{\rm hi}]$,
and the mixture is unstable at $q_{\rm mix}$; since $\Tlin$ is a supremum over unstable
temperatures, $\Tlin^{\rm mix}\ge\sqrt{q_{\rm mix}}\,\OmE/(2\pi)$ exceeds
$\Tlin^{\rm E}\le\sqrt{q_{\rm hi}}\,\OmE/(2\pi)$.
\end{proof}

\subsection{The certification method}

The chain above is deliberately small in what it assumes, and the two verifiers that
produce it share no code and no formulation.

\emph{Soundness of the cold side.} The witness is finitely supported, so the
infinite-dimensional form restricted to it is a finite sum with no bound anywhere. In the
first verifier the finite form is evaluated with every pair entry replaced by an upper
bound on opposite-sign pairs and by a lower bound on same-sign pairs, with the Coulomb
term rounded up, all at a fixed power-of-two scale and in unbounded integers where the
64-bit fast path is not provably safe; the resulting integer is an upper bound of the
scaled form, so a negative value is a certificate. The second verifier evaluates the same
statement on the \emph{unfolded} full Matsubara line $n\in\{-N,\dots,N-1\}$ with $N=400$,
retaining $a_0$, as one integer correlation $c_d=\sum_nf_nf_{n+d}$ times the exact rational
weights; it never uses the folding of Lemma~\ref{lem:laplace}, so the two verifiers agree
only if the folding is right.

\emph{Soundness of the hot side.} Positivity must hold for every vector, not only for
nonnegative ones, and this is the step where an earlier version of the second verifier was
unsound: an entrywise lower bound on off-diagonal entries implies $A\preceq\Fcal$ only on
nonnegative vectors, and at $\muS>0$ the rank-one Coulomb block produces mixed-sign
unstable directions on which that implication fails. The repaired device is a signed
enclosure: since $a_r(q)$ is decreasing in $q$, every off-diagonal entry of the head block
increases with $q$. With an integer centre $C$ and a radius $\mathcal R$ covering the two
endpoint values, each rounded outward, the direction that soundness needs is
$C_{ij}-\mathcal R_{ij}\le\Fcal_{ij}(q_{\rm lo})$ and
$\Fcal_{ij}(q_{\rm hi})\le C_{ij}+\mathcal R_{ij}$, hence
$|\Fcal_{ij}-C_{ij}|\le\mathcal R_{ij}$ on the whole window and
$\Fcal\succeq C-\mathrm{diag}(\rho)$ with $\rho_i=\lceil\frac1{h_i}\sum_j\mathcal
R_{ij}h_j\rceil$ for $h_i=2^{22}/(i+1)$ by the arithmetic-geometric inequality. The head
size $m=\max(160,2k+1)$ keeps the entire Coulomb block inside the head, so the tail is
purely phononic. Every gap of the infinite tail is at least the uniform quantity
\begin{equation}\label{eq:gap}
g=2m+1+A_m(q_{\rm hi})-a_{2m+1}(q_{\rm lo})-S_1(q_{\rm lo})-S_{2m+1}(q_{\rm lo})>0,
\end{equation}
so the tail block satisfies $T\succeq gI$; completing the square in the tail variable gives
$\Fcal\succeq C-\mathrm{diag}(\rho)-\mathrm{diag}(\theta)$ with a diagonal penalty
$\theta$ dominating $g^{-1}B^{\rm T}B$, and the remaining decision is an exact integer
matrix problem: a floating Cholesky proposes a factor, the factor is rounded to integers at
a power-of-two scale, and the exact residual must have positive diagonal and strict
diagonal dominance. Every product is guarded, and a guard that does not hold skips the
candidate rather than clamping it.

\emph{Independence.} The second verifier encloses $\pi$ from
$\pi/4=\arctan\frac12+\arctan\frac13$ while the first uses Machin's formula, so the two
$C_E$ brackets are arithmetically independent as well. Every field of the stored
certificate file is recomputed and compared: the exact Einstein bracket, the coefficient
bracket by inclusion, the box counts, and the mixture data with $\sum_ip_i=1$ and
$\sum_ip_iu_i=1$ checked exactly; $a_1$ of each mixture is pinned against a hand-computed
rational so that a swapped weight-position unpacking cannot survive.

\emph{Teeth.} Nine mutants of the two exact verifiers, each with an unmutated control
that must pass before the mutant counts, are refused with an exit code and a message
naming the statement. In the first verifier: a doubled repulsion at the cold witness of
the Einstein line; a mirrored spectrum together with the weight swap $p\to1-p$ that breaks
the second moment and the transfer of the mixture's verdict to the Einstein line at
$q_{\rm mix}$; a cold witness claimed at the stable endpoint $q_{\rm hi}$; a disabled
far-tail bound; and a sign-flipped (attractive) mask term. In the second verifier: a hot
cover extended down to $q_{\rm lo}$, where the form is indefinite; a doubled repulsion
that removes the cold sign; the weight swap refused by the second-moment guard; and a
shaved $C_E$ bracket that must not survive the receipt check. Two further self-tests have
teeth: the structural checker refuses 36 perturbed statements and its folded mode detects
the display that omits the tail term in all 120 trials, and the direct-discretization
program refuses a mirrored kernel and a shifted normal diagonal. Two regression modes
replay earlier defects: three documented boxes on which the rejected verifier returned a
positive certificate are shown to be exactly negative at a rational point and must now be
refused, and a six-window self-check recomputes the head entries in exact rational
arithmetic, 573\,003 comparisons with no floating point in the decision, where the
previous formula violated the enclosure on all 189\,950 off-diagonal entries by up to
45\,239 integer units. All twenty recorded checks pass with their expected exit codes;
the mutants refuse with the intended messages and no control broke.

\section{Certified brackets and the zero-repulsion calibration}\label{sec:brackets}

Table~\ref{tab:brackets} records the seven Einstein brackets certified in this work: the
three reference brackets of Section~\ref{sec:cert} and four further instances at or below
$\muS/\lambda=1/2$. Each row carries a cold witness at $q_{\rm lo}$, a hot certificate on
the interval shown, and a coefficient bracket. A certified Einstein bracket is a statement
about the Einstein line only. It is not a statement of optimality: only the rows of
Table~\ref{tab:cex} make an optimality claim, and they make it negatively.

\begin{table*}[t]
\centering
\footnotesize
\setlength{\tabcolsep}{5pt}
\renewcommand{\arraystretch}{1.15}
\begin{tabular}{@{}crrccccc@{}}
\toprule
 & $\lambda$ & $\muS$ & $\muS/\lambda$ & $R^2$ & Einstein $q_E$ bracket & $C_E$ bracket & hot interval\\
\midrule
E1 & $3$ & $2$ & $2/3$ & $9/4$ & $[0.0083,\,0.0095]$ & $[0.00837140,\,0.00895615]$ & $[0.0095,\,3.07270)$\\
E2 & $2$ & $3/2$ & $3/4$ & $4$ & $[0.0009,\,0.003]$ & $[0.00337618,\,0.00616405]$ & $[0.003,\,1.77000)$\\
E3 & $4$ & $4$ & $1$ & $4$ & $[129/500000,\,0.003]$ & $[0.00127820,\,0.00435864]$ & $[0.003,\,4.38414)$\\
A1 & $1$ & $1/2$ & $1/2$ & $1$ & $[0.0112,\,0.012]$ & $[0.01684337,\,0.01743456]$ & $[0.012,\,0.518654)$\\
A2 & $1$ & $1/10$ & $1/10$ & $1$ & $[0.308,\,0.3105]$ & $[0.08832740,\,0.08868516]$ & $[0.3105,\,0.518654)$\\
A3 & $3$ & $1$ & $1/3$ & $9/4$ & $[0.456,\,0.4575]$ & $[0.06205001,\,0.06215200]$ & $[0.4575,\,3.07270)$\\
A4 & $10$ & $1$ & $1/10$ & $9$ & $[4.56,\,4.565]$ & $[0.10747378,\,0.10753270]$ & $[4.565,\,12.2824)$\\
\bottomrule
\end{tabular}
\caption{The seven certified Einstein brackets. Brackets are exact rationals, outward
rounded for display; $C_E=\sqrt q/(2\pi\sqrt\lambda)$ and the hot interval is the certified
range of stability of the Einstein form, closed on the left and open on the right. The
right endpoint shown is an outward display of the exact certified upper endpoint
$q_{\rm ph}^{\rm hi}(\lambda)\ge q_{\rm ph}(\lambda)$ of the zero-repulsion bracket, beyond
which Corollary~\ref{cor:value} supplies stability, so the certified range and the
corollary together reach every larger $q$ without a gap. Box counts of the hot covers are
$19$, $36$ and $37$ for E1, E2 and E3, and $10$, $12$, $14$ and $11$ for A1, A2, A3 and
A4, in each case the second verifier's own outward count; the stored certificate file carries the first verifier's own count for each
interval, which is a different subdivision and is checked as a receipt.}
\label{tab:brackets}
\end{table*}

Two readings of the table are worth separating. Across the first three rows the coefficient
$C_E$ moves by a factor of seven and the ratio $\muS/\lambda$ by a factor of $1.5$: the
spread of the certified values across different couplings and cutoffs is not a controlled
comparison of onsets, and no monotonicity in $\muS$ is claimed from it. In the second group
$C_E$ sits strictly below the zero-repulsion value at the same $\lambda$: at $\lambda=1$
the anchor A2 has $C_E\in[0.08832740,0.08868516]$ while the zero-repulsion coefficient is
$C(1)\in[0.1146194547,0.1146195694]$ (\textsc{certified} in~\cite{finiteband,univglue}),
and at $\lambda=10$ the anchor A4 has $C_E\in[0.10747378,0.10753270]$ against
$C(10)\in[0.1763847604,0.1763849369]$. That is the expected direction, since a positive
pseudopotential can only subtract, and the two numbers are onsets of different objects (one
with repulsion through a common cutoff, one without any repulsion), so the comparison is
consistency rather than inference.

The calibration $C(1)\in[0.11461945,0.11461957]$ deserves a word, because it fixes the
absolute scale of everything above. It is the $\lambda=1$, $\muS=0$ coefficient of the
zero-repulsion ceiling~\cite{finiteband,univglue}, certified on exact rationals with
outward-rounded endpoints, and it corresponds to $\kb\Tlin=\sqrt{q}\,\OmE/(2\pi)$ with
$q\in[0.51865242539,0.51865346269]$. The chain $0.1146194547$, $0.1610715064$ at
$\lambda=3$, $0.1763847604$ at $\lambda=10$, $0.18272624777232952$ as
$\lambda\to\infty$ is monotone in $\lambda$ and converges to the certified $\gamma=2$
constant~\cite{gamma}. Nothing in the counterexamples moves it: by
Corollary~\ref{cor:value} the whole family of positive-$\muS$ onsets lies below the
$\muS=0$ Einstein value at the same $\lambda$.

\section{Direct-discretization cross-check}\label{sec:num}

A third, independent implementation builds the linearized equation directly on the full
fermionic Matsubara line with $n\in[-240,239]$, the kernel
$K_{nm}=a_{|n-m|}-\muS c_nc_m$, the closed-form normal diagonal and an explicit projector on
the even gap subspace, in float64 throughout. It shares no code with the two verifiers and
carries no certificate; its results are \textsc{numerical}. Table~\ref{tab:direct} reports
its onsets. Every Einstein onset sits inside or above its certified bracket, every
certified-unstable mixture point is unstable there, and the direct even-sector eigenvalue
agrees with the folded form to $4.4\times10^{-16}$ on all seven instances, which is the
sharpest available check that the folding of Lemma~\ref{lem:laplace} is implemented
correctly. An earlier version of this direct implementation disagreed with the folded
form by factors of $2.6$ to $7$; the cause was not the mirror term but four defects in
that version, namely bosonic instead of fermionic Matsubara frequencies, a doubled
$\lambda$, a wrong normal diagonal and an unsymmetrized kernel, and none of them survives
in the code reported here.

\begin{table}[t]
\centering
\footnotesize
\setlength{\tabcolsep}{6pt}
\begin{tabular}{@{}lccc@{}}
\toprule
case & Einstein onset $q$ (direct) & mixture onset $q$ (direct) & $\Tlin^{\rm mix}/\Tlin^{\rm E}$ (direct)\\
\midrule
E1 & $0.0083924335$ & $0.012345679$ & $1.21$\\
E2 & $0.00090451693$ & $0.0076896872$ & $2.92$\\
E3 & $0.00025806035$ & $0.0256$ & $9.96$\\
\bottomrule
\end{tabular}
\caption{Float64 onsets in the direct full-line discretization (\textsc{numerical}). Doubling
the band from $|n|\le240$ to $|n|\le480$ shifts the largest eigenvalue by at most
$6.4\times10^{-5}$ and flips no sign. The certified chain of Section~\ref{sec:cert} does not
depend on this table.}
\label{tab:direct}
\end{table}

The same code supplies the small-pseudopotential controls. With $\lambda=1$ and $R=1$ the
\emph{same} two-line spectrum as E1, evaluated at $\muS/\lambda=1/10$ and $1/2$, is stable
at the lower end of the Einstein bracket: at $q=0.308$ the Einstein form has largest
eigenvalue $1.0032$ against $0.1817$ for the mixture, and at $q=0.0112$ it has $1.0058$
against $0.4133$. In the zero-repulsion comparison, where the Einstein line is provably
optimal, the mixture is on the stable side by a wide margin. The reversal of
Section~\ref{sec:cert} is therefore a property of the repulsion and not of the mixture.

\section{Below one half: certified counterexamples, analytic regimes,
and numerical phase data}\label{sec:half}

The instances of Sections~\ref{sec:cert} and~\ref{sec:brackets} sit at
$\muS/\lambda\in\{2/3,3/4,1\}$. The natural next question is where the failure of Einstein
optimality begins, and whether $\muS/\lambda=1/2$ is a threshold. It is not, but below one
half the failure requires a coupling that is itself numerically located. This section
proceeds in three steps. Section~\ref{sec:half:cert} certifies four counterexamples
strictly below one half. Section~\ref{sec:half:proof} proves the analytic statements that
bound the failure: a two-line reduction at fixed mask, exact optimality whenever the mask
is empty at the crossing, the resulting large-coupling regime, the single-mode regime with
its sharp hypothesis, and a parameter-free positivity bound above $q=4\lambda$ that closes
the hot side of every certificate here without computation. Section~\ref{sec:half:num}
records where the certified failure stops; that part is \textsc{numerical} and carries no
certificate and no proof.

Throughout this section $p\in\mathcal P_1=\{$probability measures on $(0,\infty)$ with
$\int dp=\int u\,dp=1\}$, $\Fcal_p^{(k)}(q)$ is the form \eqref{eq:form} built from the
transfers $a_r=a_r(q;p)$ and the mask count $k=k(q,R)$ of \eqref{eq:mask}, $\mathcal D_0$
is the set of finitely supported $h$, and $q_E=q_E(\lambda,\muS,R)$ is the Einstein onset
at the instance, the value at which $\Fcal_{\delta_1}^{(k(q,R))}(q)$ changes sign. At
$\muS=0$ this onset is $q_{\rm ph}(\lambda)$ and it does not depend on $R$, in agreement
with \eqref{eq:ceiling}.

\subsection{Four certified counterexamples below one half}\label{sec:half:cert}

\begin{theorem}[Four certified failures of Einstein optimality below one half]
\label{thm:half}\certified
For each of the four parameter sets of Table~\ref{tab:half} there exists a rational
two-line spectrum with the same $(\lambda,M,\muS,\omega_c)$ as the Einstein line, hence
the same $\OmE=\sqrt{M/\lambda}$ and the same $R$, whose linear onset is strictly larger:
$\kb\Tlin^{\rm mix}\ge\sqrt{q_{\rm mix}}\,\OmE/(2\pi)$ with $q_{\rm mix}>q_{\rm hi}$,
where $[q_{\rm lo},q_{\rm hi}]$ is the certified Einstein bracket of that row. The
certified margins are $q_{\rm mix}/q_{\rm hi}\ge1.0476$, $1.0379$, $1.2171$ and $1.1893$
in the order of the table, that is temperature gains of at least $2.353$, $1.878$,
$10.326$ and $9.055$ percent. Every row has $\muS/\lambda<1/2$.
\end{theorem}

\begin{table}[t]
\centering
\scriptsize
\setlength{\tabcolsep}{2.5pt}
\renewcommand{\arraystretch}{1.15}
\begin{tabular}{@{}rrrrcccccc@{}}
\toprule
$\lambda$ & $\muS$ & $\muS/\lambda$ & $R^2$ & $q_E$ & $\CE$ & spectrum & $q_{\rm mix}$ & gain & direct\\
\midrule
$25.6$ & $12.2522$ & $0.478602$ & $11.3$ & $[0.519,\,0.525]$ &
$[0.02266124,\,0.02279187]$ & $1-6.0\times10^{-6}$ at $0.88$ & $0.55$ & $\ge2.353\%$ &
$1.03295$\\
$26$ & $12.48$ & $0.480000$ & $10.6$ & $[0.495,\,0.501]$ &
$[0.02196018,\,0.02209288]$ & $1-2.0\times10^{-5}$ at $0.8$ & $0.52$ & $\ge1.878\%$ &
$1.02446$\\
$27$ & $13.09$ & $0.484815$ & $12.6$ & $[0.414,\,0.419]$ &
$[0.01970781,\,0.01982647]$ & $1-1.55\times10^{-4}$ at $0.69$ & $0.51$ & $\ge10.326\%$ &
$1.16047$\\
$28$ & $13.692$ & $0.489000$ & $12.4$ & $[0.407,\,0.412]$ &
$[0.01918838,\,0.01930589]$ & $1-3.60\times10^{-4}$ at $0.64$ & $0.49$ & $\ge9.055\%$ &
$1.18148$\\
\bottomrule
\end{tabular}
\caption{The four counterexamples below one half, certified by two exact verifiers
(\textsc{certified}). The masses and second moments of the two-line spectra are fixed by
the moment condition, so each spectrum is written as a weight $1-\epsilon$ at $u_1$; the
small weight sits at a stiff mode $u_2\in\{20000,10000,2000,1000\}$. Einstein brackets and
coefficients have the same meaning as in Table~\ref{tab:brackets}, the gain column is the
lower bound $\sqrt{q_{\rm mix}/q_{\rm hi}}-1$ on the temperature gain, and the last column
is the float64 ratio $\Tlin^{\rm mix}/\Tlin^{\rm E}$ of the direct discretization of
Section~\ref{sec:num} (\textsc{numerical}).}
\label{tab:half}
\end{table}

\begin{proof}[Certification]
Each row is a chain of finite exact statements of the same three-part shape as in
Theorem~\ref{thm:cert}, with two differences of substance: the hot cover stops at
$q=4\lambda$, and the second verifier checks the cold side without the folding.
\emph{Cold side.} At the rational $q_{\rm lo}$ an integer witness vector supported on the
first $n$ folds is given, and the value of the infinite-model form on it is exactly the
finite quadratic form, so a negative value is a certificate with no tail bound. The first
verifier evaluates the folded form. The second evaluates the same witness on the unfolded
full Matsubara line, keeping $a_0$, as exact integer correlations times exact rational
weights, and never uses Lemma~\ref{lem:laplace}; the two agree only if the folding and the
mask are both implemented correctly.
\emph{Hot side.} On $[q_{\rm hi},4\lambda)$, split at every exact mask jump
$q=4R^2/(2\kappa-1)^2$, the Einstein form is certified strictly positive by the signed
enclosure of Section~\ref{sec:cert}, with integer scale $2^{30}$ and head $n=240$ in the
first verifier and integer scale $2^{22}$ and head $m=200$ in the second; the infinite tail
is eliminated in the first verifier by an exact rational far-tail bound, and in the second
by a per-row bound on the tail together with an exact residual and strict diagonal
dominance of the head. Above $4\lambda$ no work is needed, since
Theorem~\ref{thm:asym} certifies positivity there for every spectrum and every
$\muS\ge0$. Hence $q_{\rm lin}(\delta_1)\in[q_{\rm lo},q_{\rm hi}]$.
\emph{Comparison.} At the rational $q_{\rm mix}>q_{\rm hi}$ the same two-line spectrum has
an exact integer witness on which the folded form (first verifier) and the unfolded
full-line form (second verifier) are both strictly negative, so
$\Tlin^{\rm mix}\ge\sqrt{q_{\rm mix}}\,\OmE/(2\pi)>\sqrt{q_{\rm hi}}\,\OmE/(2\pi)\ge
\Tlin^{\rm E}$.
\emph{Margins and cost.} Every certified box carries a strictly positive exact integer
margin. The four hot covers use $2638$, $1444$, $539$ and $239$ boxes in the first
verifier and $2843$, $1504$, $541$ and $239$ in the second, and the smallest margins of
the second verifier are $0.0697854$, $0.244918$, $0.504196$ and $45.2602$ in the four
rows, quoted in that verifier's $2^{22}$-scaled integer units, each backed by a strictly
positive integer Gershgorin gap. The top row is the tight one, because its certified hot
interval begins closest to the onset, at $q_{\rm hi}=0.525$, and that cover is subdivided
finest. When the bisection depth budget is exhausted the chain raises instead of
accepting.
\emph{Teeth.} Sixteen perturbed statements across the three exact programs of this
section are refused with an exit code and a message naming the statement: a doubled
repulsion at the cold witness, a mirrored spectrum and the weight swap that breaks the
second moment, a cold witness claimed at the stable endpoint, a disabled far-tail bound, a
sign-flipped mask term, a hot cover extended down to $q_{\rm lo}$, a shaved coefficient
bracket, a dropped overflow guard, a flipped radius sign, and three structural
perturbations of Lemma~\ref{lem:decomp} and of the diagonal ordering used in
Theorem~\ref{thm:cone}; the structural checker in addition refuses when all three of its
mutants are run together. It recomputes every algebraic identity on forty random rational
two-line spectra and trial vectors.
\emph{Two refused defects.} Two defects of an earlier version of the second verifier are
now permanent refused controls, because both turned instability into false positive
certificates: an entrywise lower bound on off-diagonal entries, which implies
$A\preceq\Fcal$ only on nonnegative vectors and fails on the mixed-sign directions that a
positive $\muS$ produces, a witness being $x^{\!\top}Ax=+565.9\times10^{12}>0$ against
$x^{\!\top}\Fcal x=-85.47\times10^{12}$, both quadratic forms taken on the same
$2^{22}$ integer scale; and a tail penalty that fell short of the exact
boundary Schur term by a factor $2.207$, at $0.619931$ against a claimed $0.280929$. Three
adversarial boxes, two of which the rejected version certified positive, are exactly
negative at a rational point, with values $-8.416\times10^{12}$, $-2.898\times10^{12}$ and
$-7.708\times10^{9}$ on the integer scale, and must now be refused. A self-test
recomputes the entries of the head enclosure in exact rational arithmetic on six windows
and three interior points each, which is 720\,000 exact rational entry comparisons,
together with eighteen exact quadratic-form comparisons of the radius penalty and the
per-row tail gaps. All twenty-one recorded commands of this section return their stated
exit codes, each refusal with a message naming the statement it refuses and each base run
at zero.
\end{proof}

Theorem~\ref{thm:half} settles the sharpness of one half in the negative for this model:
$\muS/\lambda=1/2$ is not the boundary of the failure region, and no statement of the form
``Einstein is optimal for $\muS/\lambda\le1/2$'' holds across all $(\lambda,R)$.

The same two verifiers also certify four Einstein reference brackets at
$\muS/\lambda\le1/2$, collected in Table~\ref{tab:anchors}. They are brackets of the
Einstein onset, not optimality statements, and they pin the sub-one-half chain to the
zero-repulsion values of Section~\ref{sec:brackets}.

\begin{table}[t]
\centering
\footnotesize
\setlength{\tabcolsep}{5pt}
\renewcommand{\arraystretch}{1.15}
\begin{tabular}{@{}rrrcccc@{}}
\toprule
$\lambda$ & $\muS$ & $\muS/\lambda$ & $R^2$ & Einstein $q_E$ bracket & $\CE$ bracket &
hot boxes\\
\midrule
$2$ & $1$ & $1/2$ & $9/4$ & $[0.0402,\,0.0410]$ & $[0.02256410,\,0.02278752]$ & $12$\\
$3$ & $3/2$ & $1/2$ & $4$ & $[0.0593,\,0.0601]$ & $[0.02237622,\,0.02252666]$ & $18$\\
$10$ & $4$ & $2/5$ & $9$ & $[0.441,\,0.447]$ & $[0.03342253,\,0.03364914]$ & $10$\\
$20$ & $19/2$ & $0.475$ & $25$ & $[0.820,\,0.827]$ & $[0.03222641,\,0.03236368]$ & $17$\\
\bottomrule
\end{tabular}
\caption{Four certified Einstein brackets at $\muS/\lambda\le1/2$ (\textsc{certified}).
Each row carries a cold witness at $q_{E,\rm lo}$ and a strictly positive hot cover on
$[q_{E,\rm hi},4\lambda)$, whose box count is the last column, after which
Theorem~\ref{thm:asym} supplies positivity. The brackets are exact rationals, outward
rounded for display. These are reference points for the Einstein onset; they make no
optimality claim.}
\label{tab:anchors}
\end{table}

\subsection{Analytic results below one half}\label{sec:half:proof}

The four certificates stand on their own, and the five results below explain why they are
possible and where they cannot be. The first reduces every comparison at fixed mask to a
two-line spectrum, so the scan of Section~\ref{sec:half:num} is not a restriction. The
last four delimit regimes in which the repulsion provably cannot disturb the Einstein
optimum. Only Theorem~\ref{thm:nomask} uses a result from outside this paper, the
zero-repulsion strict optimality of the Einstein line~\cite{rens,kondratev}, and that
dependency is written into its hypothesis; Theorem~\ref{thm:large} applies it and inherits
the dependency.

\begin{lemma}[Exact finite-support decomposition]\label{lem:decomp}\proved
For every finitely supported $h\in\ell^2(\mathbb Z_{\ge0})$ and every $p\in\mathcal P_1$
the folded form satisfies
\begin{equation}\label{eq:finite}
\langle h,\Fcal_p^{(k)}(q)h\rangle=\sum_{i\ge0}(2i+1)h_i^2
+2\muS\Big(\sum_{i<k}h_i\Big)^2+\sum_{j\ge1}a_j(q;p)L_j(h),
\end{equation}
with
\begin{equation}\label{eq:Lj}
L_j(h)=-\sum_{i=0}^{j-1}h_ih_{j-1-i}-2\sum_{i\ge0}h_{i+j}(h_i-h_{i+j})
=\tfrac12\sum_{i=0}^{j-1}(h_i-h_{j-1-i})^2+\sum_{i\ge0}(h_i-h_{i+j})^2
-2\sum_{i=0}^{j-1}h_i^2,
\end{equation}
and $|L_j(h)|\le5\|h\|_2^2$ with $L_j(h)=0$ for $j\ge2n$ when $h$ is supported on
$\{0,\dots,n-1\}$.
\end{lemma}

\begin{proof}
By linearity in $(a_j)_{j\ge1}$, the coefficient of $a_j$ in the value of the form is
$2\sum_{i\ge j}h_i^2$ from the normal self-energy diagonal $2\sum_{m=1}^{i}a_mh_i^2$,
$-2\sum_{i\ge0}h_ih_{i+j}$ from the Toeplitz off-diagonal $-a_{|i-j|}$ at
$|i-j|=j$, and $-\sum_{i=0}^{j-1}h_ih_{j-1-i}$ from the Hankel term $-a_{i+j+1}$, whose
diagonal at $i=j$ belongs to it when $j=2i+1$. Combining the first two,
$2\sum_{i\ge j}h_i^2-2\sum_{i\ge0}h_ih_{i+j}
=\sum_{i\ge0}(h_i-h_{i+j})^2-\sum_{i=0}^{j-1}h_i^2$, and the Hankel sum rewrites as
$\tfrac12\sum_{i=0}^{j-1}(h_i-h_{j-1-i})^2-\sum_{i=0}^{j-1}h_i^2$. Adding the two gives
both expressions in \eqref{eq:Lj}. Support on $\{0,\dots,n-1\}$ makes every term in
\eqref{eq:Lj} vanish for $j\ge2n$, and \eqref{eq:finite} extends to the whole form domain
by continuity with $\sum_j a_j(q;p)<\infty$. The identity is also recomputed in exact
rational arithmetic on random spectra by the structural checker of
Section~\ref{sec:half:cert}.
\end{proof}

\begin{theorem}[Two-line reduction at fixed mask and repulsion]\label{thm:twodelta}\proved
Fix $\lambda>0$, $\muS\ge0$, $q>0$ and an integer mask count $k\ge0$, and let
$\mathcal P_1^{(2)}=\{(1-\theta)\delta_{u_1}+\theta\delta_{u_2}:0<u_1\le1\le u_2,\
\theta\in[0,1],\ (1-\theta)u_1+\theta u_2=1\}\subset\mathcal P_1$. For every
$p\in\mathcal P_1$ and every $h\in\mathcal D_0$ there is $p_2\in\mathcal P_1^{(2)}$ with
$\langle h,\Fcal_{p_2}^{(k)}(q)h\rangle=\langle h,\Fcal_p^{(k)}(q)h\rangle$, and
consequently
\[
\inf_{p\in\mathcal P_1}\lambda_{\min}\big(\Fcal_p^{(k)}(q)\big)
=\inf_{p_2\in\mathcal P_1^{(2)}}\lambda_{\min}\big(\Fcal_{p_2}^{(k)}(q)\big).
\]
\end{theorem}

\begin{proof}
Fix $h$ supported on $\{0,\dots,n-1\}$ and $q>0$. By Lemma~\ref{lem:decomp},
$L_j(h)=0$ for $j\ge2n$, so
$\langle h,\Fcal_p^{(k)}(q)h\rangle=C(h,k,\muS)+\int\Phi_{h,q}(u)\,dp(u)$ with
$\Phi_{h,q}(u)=\lambda\sum_{j=1}^{2n-1}\frac{u}{u+qj^2}L_j(h)$, a finite sum of bounded
continuous functions on $(0,\infty)$ with
$|\Phi_{h,q}(u)|\le\lambda\sum_{j=1}^{2n-1}|L_j(h)|$ uniformly, and both $u\mapsto u$ and
$\Phi_{h,q}$ are $p$-integrable for every $p\in\mathcal P_1$. By the
Richter--Rogosinski (Markov--Krein) theorem applied to the curve
$\Gamma=\{(u,\Phi_{h,q}(u)):u>0\}$, the moment point
$(1,\int\Phi_{h,q}\,dp)$ lies in $\mathrm{conv}(\Gamma)$ and is a convex combination of at
most two points of $\Gamma$; since the first moment of that combination is exactly one,
one atom satisfies $0<u_1\le1$ and the other $u_2\ge1$ whenever there are two. The
resulting $p_2$ reproduces the value of the form at $h$. Taking the infimum over
normalized $h\in\mathcal D_0$ and over $p\in\mathcal P_1$, and using
$\mathcal P_1^{(2)}\subset\mathcal P_1$, gives the two-sided equality.
\end{proof}

Theorem~\ref{thm:twodelta} is why every comparison in this paper, including the scans of
Section~\ref{sec:half:num}, may be carried out inside the two-line family: at fixed $q$,
$k$ and $\muS$ no spectrum whatever improves on the best two-line spectrum.

\begin{theorem}[Exact optimality when the mask is empty at the crossing]
\label{thm:nomask}\proved
Assume the zero-repulsion strict optimality of the Einstein line~\cite{rens,kondratev}, and
let $q_{\rm ph}(\lambda)$ be its Einstein onset. If $4R^2<q_{\rm ph}(\lambda)$, that is
$\omega_c/\OmE<\tfrac12\sqrt{q_{\rm ph}(\lambda)}=\pi C(\lambda)\sqrt\lambda$, then the
mask count vanishes on an open neighbourhood of $[q_{\rm ph}(\lambda),\infty)$ and, for
every $\muS\ge0$,
\[
q_{\rm lin}(\lambda,\muS,R;p)\le q_{\rm lin}(\lambda,\muS,R;\delta_1)
=q_{\rm ph}(\lambda),
\]
with strict inequality for every $p\in\mathcal P_1\setminus\{\delta_1\}$.
\end{theorem}

\begin{proof}
Since $4R^2<q_{\rm ph}(\lambda)$, every $q\in(4R^2,\infty)\ni q_{\rm ph}(\lambda)$
satisfies $(2i+1)^2q\ge q>4R^2$ for all $i\ge0$, so no fold is masked and
$\Fcal_p^{(k(q,R))}(q)=\Fcal_p^{(0)}(q)$ on that whole interval. That form does not involve
$\muS$, so the comparison is the zero-repulsion one: $\Fcal_p^{(0)}(q)\succ0$ for
$q\ge q_{\rm ph}(\lambda)$ and every $p\ne\delta_1$, with a strictly positive spectral gap
at $q_{\rm ph}(\lambda)$ because the operator has compact resolvent, while
$\lambda_{\min}(\Fcal_{\delta_1}^{(0)}(q_{\rm ph}(\lambda)))=0$. Hence
$q_{\rm lin}(p)\le q_{\rm ph}(\lambda)$ for all $p$, and $q_{\rm lin}(p)<q_{\rm ph}(\lambda)$
for $p\ne\delta_1$ by continuity of the spectrum in $q$ below the crossing.
\end{proof}

The strict inequality $4R^2<q_{\rm ph}(\lambda)$ cannot be relaxed with $\muS>0$: at
equality the mask contains the fold $i=0$ exactly for $q<q_{\rm ph}(\lambda)$, so the
rank-one term acts precisely where the Einstein form is indefinite and lowers the
Einstein onset.

\begin{theorem}[Large coupling at fixed cutoff ratio]\label{thm:large}\proved
For every fixed $R>0$, every $\lambda\ge4R^2+1$ and every $\muS\ge0$,
$q_{\rm ph}(\lambda)>\lambda-1\ge4R^2$ and Einstein strictly maximizes the linear onset:
\[
q_{\rm lin}(\lambda,\muS,R;p)<q_{\rm lin}(\lambda,\muS,R;\delta_1)=q_{\rm ph}(\lambda)
\qquad\forall p\in\mathcal P_1\setminus\{\delta_1\}.
\]
\end{theorem}

\begin{proof}
At $q_*=\lambda-1>0$ the first Einstein transfer is $a_1(q_*;\delta_1)=\lambda/(1+q_*)=1$,
so the $(0,0)$ entry of $\Fcal_{\delta_1}^{(0)}(q_*)$ is $1-a_1=0$, while the $(0,1)$
entry is $-(a_1+a_2)=-(1+\lambda/(1+4(\lambda-1)))<-1<0$. For
$v_\varepsilon=e_0+\varepsilon e_1$ and small $\varepsilon>0$ therefore
$\langle v_\varepsilon,\Fcal_{\delta_1}^{(0)}(q_*)v_\varepsilon\rangle
=2\varepsilon(\Fcal_{\delta_1}^{(0)}(q_*))_{01}
+\varepsilon^2(\Fcal_{\delta_1}^{(0)}(q_*))_{11}<0$, so
$\lambda_{\min}(\Fcal_{\delta_1}^{(0)}(q_*))<0$ and
$q_{\rm ph}(\lambda)>q_*=\lambda-1\ge4R^2$. Theorem~\ref{thm:nomask} then gives the
claim. The trial-vector statement is verified in exact rational arithmetic by the
structural checker of Section~\ref{sec:half:cert}.
\end{proof}

\begin{theorem}[Single-mode mask: cone monotonicity and its sharp failure]
\label{thm:cone}\proved
Fix $\lambda>0$, $\muS\ge0$ and $q$ in the single-mode window
$\tfrac49R^2\le q<4R^2$.
\begin{enumerate}[label=(\roman*),leftmargin=2.2em]
\item Let $p\in\mathcal P_1\setminus\{\delta_1\}$ and $\muS\le1+a_1(q;p)$. Then the
ground state of $\Fcal_p^{(1)}(q)$ is simple with a positive eigenvector
$h^*_0\ge h^*_1>h^*_2>\dots>0$, with $h^*_0>h^*_1$ whenever
$\muS<1+a_1(q;p)$, every $L_j(h^*)<0$, and
\begin{equation}\label{eq:conedom}
\lambda_{\min}\big(\Fcal_{\delta_1}^{(1)}(q)\big)
\le\langle h^*,\Fcal_{\delta_1}^{(1)}(q)h^*\rangle
<\langle h^*,\Fcal_p^{(1)}(q)h^*\rangle
=\lambda_{\min}\big(\Fcal_p^{(1)}(q)\big).
\end{equation}
In particular this holds unconditionally for $\muS\le1$, and on the whole branch for
$\muS/\lambda\le1/2$ with $\lambda\le2$.
\item If $\muS>1+a_1(q;p)$ the nonincreasing cone is not invariant, since
$\tilde D_{p,1}-\tilde D_{p,0}=2+2a_1(q;p)-2\muS<0$; along the soft family $u_1\to0^+$
with $\theta u_2\to1$ one has $a_1(q;p)\to0$, so the difference tends to $2-2\muS<0$ for
every $\muS>1$.
\end{enumerate}
\end{theorem}

\begin{proof}
Write $\Fcal_p^{(1)}(q)=\mathrm{diag}(\tilde D_p)-\tilde A_p$ with
$\tilde D_{p,i}=(2i+1)+a_0+2\sum_{j=1}^i a_j(q;p)+2\muS\mathbf 1_{i=0}$ and
$(\tilde A_p)_{im}=a_{|i-m|}(q;p)+a_{i+m+1}(q;p)>0$, so that $\tilde A_p$ is the
compression to the even subspace of the full-line convolution whose symbol
$\lambda\int\sum_r\frac{u}{u+qr^2}e^{ir\theta}\,dp>0$ is strictly positive by Poisson
summation of the positive Yukawa kernel; hence $\tilde A_p\succeq0$.
\emph{Cone preservation.} Extending a nonincreasing nonnegative $h$ symmetrically to the
full line and subtracting consecutive rows gives the exact identity
\begin{equation}\label{eq:telescope}
(\tilde A_ph)_i-(\tilde A_ph)_{i+1}
=\sum_{s\ge0}\big(a_s(q;p)-a_{s+1}(q;p)\big)\big(y_{i-s}-y_{i+1+s}\big)\ge0,
\end{equation}
with strict positivity somewhere for $h\ne0$, so $\tilde A_p$ preserves the cone.
\emph{Ground state.} Under $\muS\le1+a_1(q;p)$ the modified diagonal satisfies
$\tilde D_{p,0}\le\tilde D_{p,1}<\tilde D_{p,2}<\dots\to\infty$, and because
$\tilde A_p\succeq0$ we get $0<\tilde D_{p,0}-E_0\le\tilde D_{p,1}-E_0<\dots$ for
$E_0=\lambda_{\min}(\Fcal_p^{(1)}(q))$. The Birman--Schwinger operator
$K=\Delta^{-1/2}\tilde A_p\Delta^{-1/2}$ with $\Delta=\mathrm{diag}(\tilde D_{p,i}-E_0)$
is compact, self-adjoint, positive semidefinite and entrywise strictly positive, with
$\|K\|=1$, so Krein--Rutman gives a strictly positive eigenvector for the eigenvalue $1$,
hence a simple ground state with $h^*>0$ nonincreasing.
\emph{Jensen comparison.} For nonincreasing positive $h^*$ the second expression in
\eqref{eq:Lj} gives $L_j(h^*)\le-(h^*_0)^2<0$ for every $j\ge1$, while strict concavity of
$u\mapsto u/(u+qj^2)$ with $\int u\,dp=1$ gives
$a_j(q;\delta_1)=\lambda/(1+qj^2)>a_j(q;p)$ for $p\ne\delta_1$. Subtracting the two
instances of \eqref{eq:finite} cancels the kinetic and Coulomb terms and leaves
$\langle h^*,\Fcal_{\delta_1}^{(1)}(q)h^*\rangle
-\langle h^*,\Fcal_p^{(1)}(q)h^*\rangle
=\sum_j(a_j(q;\delta_1)-a_j(q;p))L_j(h^*)<0$, which is \eqref{eq:conedom}. Part (ii) is
the displayed diagonal difference together with the limit $a_1(q;p)\to0$ along the soft
family. The identity \eqref{eq:telescope} and the diagonal orderings over their full
proved ranges are recomputed in exact rational arithmetic by the structural checker, which
also carries explicit refutation witnesses outside the proved range.
\end{proof}

\begin{theorem}[Parameter-free positivity above $q=4\lambda$]\label{thm:asym}\proved
For every $\lambda>0$, $\muS\ge0$, $R>0$, $p\in\mathcal P_1$ and $q\ge4\lambda$ the
folded form is strictly positive definite, with
\begin{equation}\label{eq:asymstrict}
\langle h,\Fcal_p^{(k(q,R))}(q)h\rangle
\ge\Big(1-\frac{\pi^2}{12}\Big)\sum_{i\ge0}h_i^2>\frac16\|h\|_2^2
\qquad\forall h\in\ell^2(\mathbb Z_{\ge0}).
\end{equation}
\end{theorem}

\begin{proof}
Since $\muS\ge0$ the rank-one term in \eqref{eq:finite} is nonnegative, so it suffices to
take $\muS=0$. In row $i$ the sum of the absolute values of the off-diagonal entries is
$2\sum_{j=1}^i a_j+2\sum_{j>i}a_j-a_{2i+1}$, because the Toeplitz and Hankel entries
combine as $\sum_{m<i}a_{i-m}+\sum_{m>i}a_{m-i}+\sum_{m\ge0}a_{i+m+1}-a_{2i+1}$.
Subtracting this from the diagonal $d_i+2A_i$ cancels $2\sum_{j\le i}a_j-a_{2i+1}$ and
leaves the exact row margin $d_i-2\sum_{j>i}a_j\ge(2i+1)-2\sum_{j\ge1}a_j$. Since
$a_j(q;p)\le\lambda/(qj^2)$ by $\int u\,dp=1$, at $q\ge4\lambda$
\[
2\sum_{j\ge1}a_j(q;p)\le\frac{2\lambda}{q}\cdot\frac{\pi^2}{6}
\le\frac{\pi^2}{12}<0.8225 .
\]
By Schur's test, $2|h_ih_m|\le h_i^2+h_m^2$, so
$\langle h,\Fcal_p^{(0)}(q)h\rangle\ge\sum_i(2i+1-\pi^2/12)h_i^2
\ge(1-\pi^2/12)\|h\|_2^2$. The exact row margins are recomputed in exact rational
arithmetic by the structural checker.
\end{proof}

Theorem~\ref{thm:asym} is why the certificates of this section and of
Section~\ref{sec:cert} need a hot cover only on a bounded interval, and it holds for every
spectrum, every mask and every $\muS$ without computation. The analytic regimes do not
meet: Theorem~\ref{thm:large} covers couplings above $4R^2+1$,
Theorem~\ref{thm:cone} covers $\muS\le1$ or $\lambda\le2$ with $\muS/\lambda\le1/2$, and
Theorem~\ref{thm:nomask} covers cutoffs below the crossing. The window of the certified
counterexamples lies outside all three, and it also lies outside every physically
realistic coupling.

\subsection{Where the certified failure stops: numerical phase data}\label{sec:half:num}

What remains is the interior, and there the evidence is floating point only. Let
\[
r_*(\lambda,R)=\inf\big\{\muS/\lambda:\text{some spectrum beats Einstein at }
(\lambda,R)\big\},
\]
an infimum over two-line spectra, which by Theorem~\ref{thm:twodelta} loses nothing at
fixed mask. Table~\ref{tab:phase} reports $r_*$ on a grid of couplings and cutoff ratios.
Three features are read off it.

\begin{enumerate}[label=(\arabic*),leftmargin=2.2em]
\item No entry breaks one half for $\lambda<25.55$ and no entry breaks it with mask count
$k\in\{0,1,2,3,4\}$. Two kinds of cell need reading with care. The cells marked $\infty$
record that no two-line spectrum was found to beat the Einstein line at any scanned
pseudopotential, which is weaker than optimality; where the cutoff also sits below the
crossing point, Theorem~\ref{thm:nomask} upgrades that to a proof. The rows whose infimum
reads one half are bounded by the smallest ratio on the grid, so they say that the
threshold never moves below one half there, not that it touches one half with a margin.
For $\lambda\le2$ the value one half is in addition covered on the single-mode branch by
Theorem~\ref{thm:cone}, since $\muS/\lambda\le1/2$ forces $\muS\le1$.
\item Above a critical coupling $\lambda\approx25.55$ a narrow re-entrant window opens,
$\lambda\in[25.55,34.0]$ with $\omega_c/\OmE\in[3.2,3.6]$ and mask count $k=5$, in which
$r_*$ dips strictly below one half, to $r_{\min}\approx0.4785$ at
$\lambda\approx25.58$ and $\omega_c/\OmE\approx3.36$. The four certified rows of Table~\ref{tab:half} sit in this
window, at $\lambda\in\{25.6,26,27,28\}$ and $\muS/\lambda\in[0.478602,0.489]$. The
mechanism is an extreme form of the warm interval of Section~\ref{sec:cert}: at
$\lambda=26$ the $k=5$ Coulomb penalty lifts the Einstein eigenvalue above zero on a narrow
island $q\in[0.501,0.5234]$ below the mask jump at $q=4R^2/81=0.523457$, while the
mixture is unstable at $q=0.52$ inside that island; the jump to $k=4$ then truncates the
upper unstable branch of the Einstein line, leaving it stable for all $q\ge0.501$.
\item On the single-mode branch the analytic regime of Theorem~\ref{thm:cone} is not the
whole story. At $\muS/\lambda=1/2$ one already has $h^*_0<h^*_1$ for the Einstein line at
$\lambda=20$, $\muS=10$ and $q=1.86444$; Einstein optimality nevertheless holds numerically
across the two-line simplex for every $\lambda\in(2,151]$, and at $\lambda=152$ with
$\muS=76$ the two-line spectrum with $(u_1,u_2)=(10^{-5},17.5)$ has
$\lambda_{\min}(\Fcal_p^{(1)}(q_E))=-6.717\times10^{-5}<0$ at $q_E=4.33295$. No certificate
is claimed for that failure either.
\end{enumerate}

\begin{table}[t]
\centering
\scriptsize
\setlength{\tabcolsep}{2.6pt}
\renewcommand{\arraystretch}{1.1}
\begin{tabular}{@{}rrrrrrrrrrrr@{}}
\toprule
$\lambda\backslash\omega_c/\OmE$ & $1$ & $1.5$ & $2$ & $2.5$ & $3$ & $3.36$ & $3.5$ & $4$ & $5$ &
$\inf_R$ & $k$\\
\midrule
$0.5$ & $\infty$ & $0.578$ & $0.685$ & $0.768$ & $0.835$ & $0.875$ & $0.890$ & $0.941$ &
$1.031$ & $0.5000$ & $1$\\
$1$ & $0.518$ & $0.608$ & $0.692$ & $0.763$ & $0.822$ & $0.858$ & $0.871$ & $0.916$ &
$0.992$ & $0.5000$ & $1$\\
$2$ & $0.500$ & $0.562$ & $0.625$ & $0.680$ & $0.728$ & $0.758$ & $0.770$ & $0.808$ &
$0.874$ & $0.5000$ & $1$\\
$3$ & $\infty$ & $0.525$ & $0.578$ & $0.626$ & $0.668$ & $0.695$ & $0.705$ & $0.739$ &
$0.799$ & $0.5000$ & $1$\\
$5$ & $\infty$ & $0.500$ & $0.528$ & $0.565$ & $0.600$ & $0.622$ & $0.630$ & $0.660$ &
$0.712$ & $0.5000$ & $1$\\
$10$ & $\infty$ & $\infty$ & $\infty$ & $0.511$ & $0.533$ & $0.548$ & $0.554$ & $0.576$ &
$0.616$ & $0.5000$ & $1$\\
$20$ & $\infty$ & $\infty$ & $\infty$ & $\infty$ & $0.500$ & $0.508$ & $0.511$ & $0.525$ &
$0.552$ & $0.5000$ & $1$\\
$25.6$ & $\infty$ & $\infty$ & $\infty$ & $\infty$ & $\infty$ & $0.4785$ & $0.502$ &
$0.512$ & $0.536$ & $0.4785$ & $5$\\
$26$ & $\infty$ & $\infty$ & $\infty$ & $\infty$ & $\infty$ & $0.4795$ & $0.501$ &
$0.511$ & $0.534$ & $0.4795$ & $5$\\
$27$ & $\infty$ & $\infty$ & $\infty$ & $\infty$ & $\infty$ & $0.4835$ & $0.4828$ &
$0.508$ & $0.530$ & $0.4828$ & $5$\\
$28$ & $\infty$ & $\infty$ & $\infty$ & $\infty$ & $\infty$ & $0.4875$ & $0.4860$ &
$0.506$ & $0.527$ & $0.4860$ & $5$\\
$30$ & $\infty$ & $\infty$ & $\infty$ & $\infty$ & $\infty$ & $\infty$ & $0.4910$ &
$0.502$ & $0.521$ & $0.4910$ & $5$\\
$35$ & $\infty$ & $\infty$ & $\infty$ & $\infty$ & $\infty$ & $\infty$ & $0.500$ &
$0.504$ & $0.512$ & $0.5000$ & $1$\\
$50$ & $\infty$ & $\infty$ & $\infty$ & $\infty$ & $\infty$ & $\infty$ & $\infty$ &
$\infty$ & $0.504$ & $0.5000$ & $1$\\
\bottomrule
\end{tabular}
\caption{The threshold ratio $r_*(\lambda,\omega_c/\OmE)$ of two-line spectra against the
Einstein line, float64 (\textsc{numerical}). The entry $\infty$ marks parameter sets at
which no two-line spectrum beats the Einstein line at any scanned pseudopotential; such a
cell is unbeaten on the grid, which is weaker than optimality. The penultimate column takes
the infimum over cutoffs of the smallest scanned pseudopotential ratio at which a two-line
spectrum beats the Einstein line, on a grid that contains one half, and $k$ is the mask
count recorded at that cutoff. These numbers locate the failure; they do not bound it.}
\label{tab:phase}
\end{table}

Two further numbers connect the certified rows to floating point. Running the direct
discretization of Section~\ref{sec:num} on folded windows of $300$ and $600$ Matsubara
folds, and on the unfolded full line, gives $\Tlin^{\rm mix}/\Tlin^{\rm E}=1.03295$,
$1.02446$, $1.16047$ and $1.18148$ for the four rows of Table~\ref{tab:half}, stable to
seven significant figures under the window doubling, and every one of them exceeds the
certified lower bound $1.0235$, $1.0187$, $1.1032$ and $1.0905$ of the same row. The
float64 ratios and the exact certificates therefore agree in direction and in magnitude.
The four Einstein anchors of Table~\ref{tab:anchors} lie below the zero-repulsion ceiling
at the same coupling, as they must, since a positive pseudopotential can only subtract
from the pairing scale.

The certified content of this section is therefore the following and nothing more: at
$\lambda\in\{25.6,26,27,28\}$ with the cutoffs and small pseudopotentials of
Table~\ref{tab:half}, Einstein is beaten; at the four parameter sets of
Table~\ref{tab:anchors}, the Einstein onset lies inside the stated bracket; and
Theorems~\ref{thm:twodelta} to~\ref{thm:asym} hold as stated. Everything about the shape of
$r_*$, about the critical coupling, about $r_{\min}$ and about the single-mode branch
beyond $\lambda=2$ is \textsc{numerical}.

\section{What is not claimed}\label{sec:not}

\begin{enumerate}[label=(\arabic*),leftmargin=2.2em]
\item The statements concern the truncated pairing-only model with a strict common cutoff
in the linearized, isotropic, flat-density flat-band setting. They are not statements
about a material, about an exact many-body transition, or about the full Coulomb
interaction. $\Tlin$ is a supremum over unstable temperatures of the linearized form, not
a root of a nonlinear gap equation and not a measured $\kb T_c$.
\item The counterexamples are kernel-level statements. The three certified ones of
Section~\ref{sec:cert} have $\muS\in\{1.5,2,4\}$ and $\muS/\lambda\ge2/3$, and the four of
Section~\ref{sec:half} have $\muS\approx12.3$ to $13.7$ with $\muS/\lambda\approx0.48$ at
$\lambda=25.6$ to $28$, where the mixture places a weight of order $10^{-6}$ to $10^{-4}$
at a mode $32$ to $141$ times $\OmE$ and the mechanism relies on the sharp step of the
mask at a Matsubara fold. Both properties put those four instances far from any physical
regime, and that is why they are reported as a statement about the model and not as a
physical anomaly. The conventional range for
phonon-mediated superconductors is $\muS\approx0.1$ to $0.15$ with $\lambda\gg\muS$. The
renormalization of $\muS$ contains no $\lambda$ and imposes no bound of the form
$\muS\le\lambda/2$, so a large kernel ratio is well defined in this model for any
$\omega_c>0$; whether such parameters are reachable in a material is not addressed. No
material prediction is made anywhere in this paper.
\item No claim is made that the ceiling value fails. Corollary~\ref{cor:value} keeps
$\kb\Tlin\le C(\lambda)\sqrt M$ at every $\muS\ge0$. What fails is the assertion that the
Einstein line attains or defines the optimum at fixed $(\lambda,M)$.
\item No claim of priority is made. The zero-repulsion ceiling, its sharp constant and the
strict single-line optimality are prior work~\cite{rens,kondratev,finiteband,gamma}; the
contribution here is the positive-pseudopotential common-cutoff question and its
certified answers above and below one half.
\item No claim is made for other Coulomb kernels. A repulsion that also enters $Z$, a
subtracted kernel, or a frequency-dependent pseudopotential is outside this model and can
change the rank-one structure of Proposition~\ref{prop:rank}, which is what both
Corollary~\ref{cor:value} and the mechanism of Remark~\ref{rem:mech} rest on.
\item Section~\ref{sec:half} is split. Its four counterexamples and four Einstein anchors
are certified by its own pair of exact verifiers, in the same two roles as
Section~\ref{sec:cert}, and its five analytic statements are proved, with the one
dependency of Theorem~\ref{thm:nomask} on the zero-repulsion theorem written into its
hypothesis and inherited by Theorem~\ref{thm:large}. The threshold ratio $r_*$, the
critical coupling $\lambda\approx25.55$, the single-mode window beyond $\lambda=2$ and the
phase table are \textsc{numerical}: they locate the failure and do not bound it. Neither
the exact threshold curve nor optimality for $\muS/\lambda\le1/2$ in general is settled
here.
\end{enumerate}

\section{Open problems}\label{sec:open}

\begin{enumerate}[label=(\arabic*),leftmargin=2.2em]
\item \emph{Einstein optimality below one half.} This is now a question about a bounded
region rather than the whole sector. Einstein is optimal when the mask dies at the
crossing (Theorem~\ref{thm:nomask}), for $\lambda\ge4R^2+1$
(Theorem~\ref{thm:large}), and on the single-mode branch for $\muS\le1$ or for $\lambda\le2$
with $\muS/\lambda\le1/2$ (Theorem~\ref{thm:cone}); and it fails, certified, in the
$k=5$ window of Section~\ref{sec:half:cert}. In every configuration tested with
$\muS/\lambda\le1/2$,
$\lambda<25.55$ and any mask width, it is not beaten. What is missing is a comparison of
forms, not of transfers: at $\muS>0$ the rank-one penalty must be shown to cost Einstein no
more than it costs an arbitrary spectrum throughout the masked region, at every mask width
$k$ a common cutoff can produce. The two-line reduction of
Theorem~\ref{thm:twodelta} reduces that to the two-line family at fixed mask, and the
interior between the proved regimes is where a proof would have to work. No argument in
either direction is available there.
\item \emph{The exact threshold curve.} Determine
$r_*(\lambda,R)=\inf\{\muS/\lambda:\text{some spectrum beats Einstein}\}$ as a function of
$(\lambda,R)$ and decide whether it has a positive lower bound. The dip below one half is
certified at four parameter sets and located numerically on a grid; its full extent, the
critical coupling of the $k=5$ window, and the failure of Einstein optimality on the
single-mode branch from $\lambda=152$ are numerical. A certified answer to any of these
would close part of the question.
\item \emph{From optimizer to value.} Corollary~\ref{cor:value} gives
$\kb\Tlin\le C(\lambda)\sqrt M$ at every $\muS\ge0$, but it does not say what the supremum
over spectra is at positive $\muS$, nor whether it is attained by any spectrum. The
degeneracy found for a spectrum-dependent cutoff shows that the supremum can equal the
zero-repulsion Einstein value in that convention. At a common cutoff the question is open,
and a natural sharper form of it is whether the supremum at $\muS>0$ is attained by a
two-line spectrum and whether it is strictly below $C(\lambda)\sqrt M$.
\item \emph{Robustness of the rank-one structure.} Which Coulomb kernels keep the ceiling
value? Proposition~\ref{prop:rank} and Corollary~\ref{cor:value} use only positivity of the
subtraction, so any pairing-only kernel that is positive semidefinite in the folded
stability form keeps $\kb\Tlin\le C(\lambda)\sqrt M$. A kernel that also modifies $Z$ does
not, and the common-cutoff optimality question for that class is not posed here.
\item \emph{Nonlinear and multichannel extensions.} The linearized onset is one step away
from the instability of the nonlinear gap equation, and the spin-chain representation of
Migdal--Eliashberg theory offers a route to that step. Extending the present comparison to
coupled channels with distinct density-of-states profiles, in the setting
of~\cite{univglue,optimum,anisosharp,capstone}, would require the same rank-one estimate
channel by channel, and nothing here settles it.
\end{enumerate}

\section*{Acknowledgments and disclosures}

AI agents generated the proofs, code and drafts under D. Rens' direction, and separate AI
agents carried out the independent reviews.

This preprint has not been peer reviewed.

\section*{Data availability}

Every constant quoted in this paper is reproduced by a replay supplement that accompanies
the preprint. The supplement contains, for the certified chain, two independently written
verifiers in exact rational and exact integer arithmetic, the fixture list of exact
rational spectra with their Einstein brackets and coefficient brackets, the stored
certificate records with their box counts and receipts, the nine mutant controls with
their expected refusals, and the two regression modes described in
Section~\ref{sec:cert}. The chain of Section~\ref{sec:half} is reproduced by its own pair
of verifiers in the same two roles, with an exact structural checker beside them: the four
sub-one-half counterexamples and the
four Einstein anchors are certified in exact arithmetic on both the folded and the unfolded
form, the five analytic statements are checked structurally with their refutation witnesses
outside the proved ranges, and sixteen mutant controls with their expected refusals
accompany them. The threshold curve and the phase table of Section~\ref{sec:half:num} come
from a separate float64 scan and are labelled as numerical. The
direct-discretization cross-check of Section~\ref{sec:num} is a float64 program and is part
of the supplement as a numerical sanity check; it is not part of any proof. No external
data set is used anywhere in this work, and no measured transition temperature is
introduced at any point.

\bibliographystyle{plain}
\bibliography{refs}
\end{document}
